% STATUS: verified
\chapter{The data model}\label{ch:data}

Everything the system knows lives in one PostgreSQL database with the \code{pgvector}
extension: the legal corpus and its vectors, the register of obligations, the companies and
their documents, the reports, the work products, the prompts, the queue of jobs, the record
of every model call and the audit trail. Only the API (Django) connects to it; the assistant
service reads and writes through the API (\cref{ch:arch}). This chapter describes the tables
as the model files define them (\file{<app>/models.py} of each application of the API). It
opens with the conventions every table follows and two figures of the relations, then goes
through the tables application by application with the columns that matter and the shape of
every JSON column, and closes with what a report copies so that it reads the same later, and
with who --- a model, the code or a person --- writes which columns. What the pipelines do with
the tables belongs to the chapters on each pipeline.

\section{Conventions}\label{sec:data:conventions}

Every application of the API follows the same rules for its tables; \file{users/models.py}
states them in its banner and the others keep them.

\paragraph{Keys.} A primary key is a \code{TextField} holding a version-4 UUID as text, made in
Python at the moment of the write (\code{id=str(uuid.uuid4())}); the database has no sequence
and assigns nothing. Because the id exists before the row is written, a step that creates
hundreds of rows writes them with one \code{bulk_create} and still holds every id --- the
gap analysis writes its findings this way, and the action plan then lists those ids in
\code{finding_ids}. Two tables are keyed by a natural key instead: \code{rate_limit_windows}
by the limiter's key and \code{model_prices} by the model's id. One table takes its ids from
outside: a row of \code{chat_messages} keeps the id the chat client gave the message, which
\file{chat/store.py} accepts only if it matches \code{MESSAGE_ID} (letters, digits and
\code{_ . : -}, at most 128 characters) and does not already belong to another thread.

\paragraph{Time.} A moment is a \code{DateTimeField} holding an aware UTC value
(\file{settings.py}: \code{USE_TZ = True}, \code{TIME_ZONE = "UTC"}). No column fills itself
in: there is no \code{auto_now}, and every write passes the value of \code{utc_now()}
(\file{common/timestamps.py}), which its banner calls \enquote{the only clock the code reads}.
On the wire a moment is an ISO~8601 string ending in \code{+00:00}. A day without a time ---
the date a version takes effect, the date an act applies from, an action's due date, the
validity of a credential --- is a \code{DateField}.

\paragraph{Enumerations.} The allowed values of a status or a kind are a tuple of constants in
the model file (\code{REPORT_STATUSES}, \code{FINDING_STATUSES}, \dots), and a
\code{CheckConstraint} in the model's \code{Meta.constraints} makes the database refuse any
other value. The constraints are named after the table and what they check (\code{reports_status_check}). The API's views check a value
against the same tuple before they write it --- \code{DOCUMENT_KINDS} when a document is
uploaded, for instance --- so the constraint guards against a value the code itself gets wrong.


A few text columns carry no constraint on purpose, and the reader meets them below:
\code{work_products.kind} (\enquote{a kind added later needs a module, not a migration}),
\code{report_acts.doc_type} (it holds the catalogue's types and also \code{internal}, the row
a company's own requirements hang on), the \code{category} of findings and actions (constrained
only on \code{obligations}), \code{report_findings.source_kind} and
\code{report_findings.override_status} (the view checks it against \code{FINDING_STATUSES}),
\code{prompt_versions.source}, \code{chat_messages.role} and the \code{outcome} and
\code{purpose} of \code{llm_calls}. Two columns of
\code{report_citations} are either empty or one of a list: their constraints allow
\code{NULL} explicitly.

\paragraph{JSON columns.} Structured values that are read and written whole --- a reading of
an act's sanctions, a finding's evidence, a report's statistics --- are \code{JSONField}
columns, stored as \code{jsonb}. Their shape is set by the code that writes them, not by the
database, and the keys inside them are mostly in camel case (\code{maxFineEur},
\code{promptVersionId}) as the API serves them --- an enumeration value used as a key keeps its
own spelling (\code{not_applicable} in \code{stats.findings}) --- while the columns themselves are
in snake case. The readers are
tolerant by design: \code{inventory_citations} in \file{companies/inventory.py}, for instance,
reads a missing or malformed list as no citations, never as an error.
\Cref{sec:data:corpus,sec:data:companies,sec:data:analysis,sec:data:programme,sec:data:ops}
give the shape of each.

\paragraph{Vectors and full text.} An embedding is a \code{VectorField} of 1,536 dimensions
(\code{EMBED_DIMENSIONS}, declared in \file{corpus/models.py} and \file{companies/models.py}
as \enquote{one model}). The model is set in \file{settings.py} (\code{AI_EMBED_MODEL}, default
\code{gateway:azure/text-embedding-3-large}), and the assistant service asks it for that
number of dimensions (\code{EMBED_DIMENSIONS} in its \file{routes/internal.js}); a change of
model means a re-embedding job, not a change of the schema (\cref{ch:models}). Four columns
hold vectors. The two that are searched --- \code{legal_chunks.embedding} and
\code{company_document_chunks.embedding} --- each have an HNSW index for cosine distance
(\code{m=16}, \code{ef_construction=64}, operator class \code{vector_cosine_ops}) and, beside
it, a GIN index over \code{to_tsvector('simple', text)} for the keyword half of the hybrid
search. The \code{simple} configuration is chosen because, in the words of
\file{corpus/models.py}, \enquote{Lithuanian has no Postgres stemmer, the vector half carries
the inflection}. The two vectors of an obligation (\code{embedding}, \code{embedding_en}) have
no index: they are the query, compared with the indexed company passages, never what is
searched.

\paragraph{What happens on delete.} Each foreign key names what happens to the row when the row
it points at is deleted, and the choices decide what survives a deletion
(\cref{tab:data:ondelete}). A part that has no meaning without its whole goes with it
(\code{CASCADE}); a reference kept for navigation is emptied and the row survives
(\code{SET_NULL}); an owner cannot be deleted while it owns something (\code{PROTECT}); and a
record that must outlive what it names holds the id as plain text, with no foreign key at all.

\begin{xltabular}{\textwidth}{@{}L{22mm}YL{44mm}@{}}
\caption{The four delete rules and where each is used}\label{tab:data:ondelete}\\
\toprule
\hd{Rule} & \hd{Where} & \hd{Consequence} \\
\midrule
\endfirsthead
\toprule
\hd{Rule} & \hd{Where} & \hd{Consequence} \\
\midrule
\endhead
\bottomrule
\endlastfoot
\code{CASCADE} & An act's versions, units and chunks, and its obligations; a version's units and chunks; a unit's children; a company's credentials, adopted standards, roles, reports, action register, work products and own requirements (\code{obligations.company}); a document's chunks; a report's acts, findings, actions and citation rows; a job's steps; a prompt's versions; a thread's messages; a user's sessions and chat threads; an adoption when the standard's act is deleted. & Deleting an act from the corpus removes its whole register of obligations; deleting a company removes every report made of it. \\
\code{SET_NULL} & Every column naming a user who acted (\code{created_by}, \code{reviewed_by}, \code{approved_by}, \dots, and \code{jobs.created_by}, \code{llm_calls.user}, \code{audit_log.user}); every link from a report's rows to the corpus and to the company's documents; an obligation's version, provision and source document; \code{legal_documents.current_version}, \code{prompts.active_version}; \code{company_documents.company}; \code{reports.job}, \code{work_products.job}; a finding's act row (\code{report_findings.act}); the company and role of a chat thread; the evidence document of a credential and the certificate of an adoption; the reports an action of the register names. & A report survives the deletion of an act, an obligation or a document it read: it holds its own copies (\cref{sec:data:snapshots}) and loses only the link. \\
\code{PROTECT} & \code{companies.owner}, \code{company_documents.owner}. & A user who owns a company or a document cannot be deleted while it does. The plain delete of an account deactivates it (\code{users.active}), which keeps the audit rows and the ownership intact; only an administrator's purge (\route{DELETE}{/api/users/{id}?purge=true}, \code{purge_account} in \file{users/purge.py}) deletes an account for good, removing its documents first, then its companies, then the account. \\
No foreign key & \code{job_steps.prompt_version_id}; \code{llm_calls.thread_id}, \code{.job_step_id}, \code{.prompt_version_id}; \code{obligations.prompt_version_id}; \code{chat_messages.parent_id}; \code{audit_log.target_id}; \code{company_actions.plan_action_id}; \code{report_actions.finding_ids} and every id inside a JSON column. & The record of who called which model, at what cost, under which prompt version, stays whole when the thread, the job step or the version it names is gone. \\
\end{xltabular}

Two deletions do more than the rules alone. Deleting a company first detaches its documents
(\code{company_documents.company} set empty, in \file{companies/api/company_views.py}): the
files belong to their owner, not to the company. Deleting a company document leaves the
requirements read from it (\code{obligations} of source \code{internal}) with an empty
\code{company_document}, and a report no longer reads them, since the internal step takes only
the requirements of the documents it reads (\file{analysis/internal.py}). Old rows of a few
tables --- expired sessions, finished jobs, usage records --- are purged by the daily
maintenance job (\cref{ch:arch,ch:ops}).

\section{The relations at a glance}\label{sec:data:figures}

\Cref{fig:data:law} shows the tables of the law, of the company and of the reports, and how they
hang together; \cref{fig:data:ops} the tables that run the system. A solid line is a part that
is deleted with its whole, a dotted line a reference that is emptied when its target goes, and
a dashed line an id held as text with no foreign key. The figures leave the columns out; the
sections after them give those.

\begin{figure}[H]
\centering
\begin{tikzpicture}[x=1mm, y=1mm, font=\sffamily\scriptsize,
    tbl/.style={svc, minimum height=7mm, inner ysep=3pt, inner xsep=4pt},
    own/.style={draw=ink, line width=.7pt},
    ref/.style={draw=muted, line width=.6pt, densely dotted}]
  \begin{scope}[on background layer]
    \draw[lane] (0,8) rectangle (68,-54);
    \draw[lane] (76,8) rectangle (150,-48);
    \draw[lane] (0,-60) rectangle (150,-114);
  \end{scope}
  \node[lanelabel, anchor=south west] at (1,-53) {the law};
  \node[lanelabel, anchor=south west] at (77,-47) {the company};
  \node[lanelabel, anchor=south west] at (1,-113) {reports and what follows them};

  % the law
  \node[tbl] (ld) at (22,-4) {\code{legal_documents}};
  \node[tbl] (lv) at (22,-18) {\code{legal_document_versions}};
  \node[tbl] (lu) at (22,-32) {\code{legal_units}};
  \node[tbl] (lc) at (22,-46) {\code{legal_chunks}};
  \node[tbl] (ob) at (54,-30) {\code{obligations}};

  % the company
  \node[tbl] (co) at (134,-22) {\code{companies}};
  \node[tbl, align=center] (cx) at (132,-1) {\code{company_credentials}\\ \code{company_standards}\\ \code{company_roles}};
  \node[tbl] (cd) at (98,-22) {\code{company_documents}};
  \node[tbl] (cc) at (100,-38) {\code{company_document_chunks}};

  % reports
  \node[tbl] (re) at (80,-70) {\code{reports}};
  \node[tbl] (ra) at (22,-80) {\code{report_acts}};
  \node[tbl] (rf) at (52,-102) {\code{report_findings}};
  \node[tbl] (rc) at (106,-86) {\code{report_citations}};
  \node[tbl] (rx) at (100,-106) {\code{report_actions}};
  \node[tbl] (ca) at (128,-70) {\code{company_actions}};
  \node[tbl] (wp) at (128,-104) {\code{work_products}};

  % the law
  \draw[own] (ld) -- node[lbl, right] {versions} (lv);
  \draw[ref] ([xshift=-6mm]ld.south) -- node[lbl, left] {current} ([xshift=-6mm]ld.south |- lv.north);
  \draw[own] (lv) -- node[lbl] {units} (lu);
  \draw[own] (lu) -- node[lbl] {chunks} (lc);
  \draw[own] ([yshift=1.2mm]lu.west) -- ++(-3,0) |- node[lbl, pos=.25, left] {parent} ([yshift=-1.2mm]lu.west);
  \draw[own] (ld.east) -| node[lbl, pos=.3] {act} (ob.north);
  \draw[ref] (lu.east) -- node[lbl, pos=.5, below] {provision} (ob.west);

  % the company
  \draw[own] (co) -- node[lbl] {documents} (cd);
  \draw[own] (cd) -- node[lbl] {chunks} (cc);
  \draw[own] (co) -- (cx);
  \draw[ref] (ob.east) -- node[lbl, pos=.45] {read from} (cd.west);
  \draw[own] (cx.north) -- ++(0,4) -| node[lbl, pos=.25] {a standard the company adopted} (ld.north);

  % reports
  \draw[own] (co.south west) -- ++(0,-28) -| node[lbl, pos=.2] {reports} (re.north);
  \draw[own] (co.south) -- node[lbl, pos=.55] {register} (co.south |- ca.north);
  \draw[own] (co.east) -- ++(4,0) |- node[lbl, pos=.3] {products} (wp.east);
  \draw[own] (re) -- node[lbl, pos=.4] {acts} (ra);
  \draw[own] (re) -- node[lbl, pos=.45] {findings} (rf);
  \draw[own] (re.south east) -- node[lbl, pos=.55, right] {citation audit} (rc.north west);
  \draw[own] (re.south) -- node[lbl, pos=.6] {plan} (rx.north west);
  \draw[ref] (ra) -- node[lbl] {act} (rf);
  \draw[ref] (rf.north) -- node[lbl, pos=.62] {obligation} (ob.south);
  \draw[ref] (ca.west) -- node[lbl, pos=.4] {report} (ca.west -| re.east);
\end{tikzpicture}
\caption{The law side and the company side. Solid: a part deleted with its whole; dotted: a
  reference emptied when its target is deleted. \code{report_actions} name their findings by id
  in a JSON list (\code{finding_ids}), not by a foreign key.}\label{fig:data:law}
\end{figure}

\begin{figure}[tbp]
\centering
\begin{tikzpicture}[x=1mm, y=1mm, font=\sffamily\scriptsize,
    tbl/.style={svc, minimum height=7mm, inner ysep=3pt, inner xsep=4pt},
    own/.style={draw=ink, line width=.7pt},
    ref/.style={draw=muted, line width=.6pt, densely dotted},
    txt/.style={draw=mid, line width=.6pt, dashed}]
  \node[tbl] (us) at (75,0) {\code{users}};
  \node[tbl] (se) at (20,0) {\code{sessions}};
  \node[tbl] (au) at (130,0) {\code{audit_log}};
  \node[tbl] (ct) at (26,-20) {\code{chat_threads}};
  \node[tbl] (cm) at (40,-40) {\code{chat_messages}};
  \node[tbl] (jo) at (75,-20) {\code{jobs}};
  \node[tbl] (js) at (75,-40) {\code{job_steps}};
  \node[tbl] (pr) at (130,-20) {\code{prompts}};
  \node[tbl] (pv) at (130,-40) {\code{prompt_versions}};
  \node[tbl] (lc) at (75,-60) {\code{llm_calls}};
  \node[tbl] (mp) at (130,-60) {\code{model_prices}};
  \node[ext] (cx) at (20,-60) {companies, roles};
  \node[ext] (rw) at (51,-20) {reports, work products};
  \node[tbl, align=center] (st) at (124,-76) {\code{mcp_servers}\quad\code{rate_limit_windows}};

  \draw[own] (us) -- node[lbl] {sessions} (se);
  \draw[ref] (us) -- node[lbl] {user} (au);
  \draw[own] (us.south west) -- node[lbl, pos=.45] {threads} (ct.north east);
  \draw[ref] (us) -- node[lbl] {created\_by} (jo);
  \draw[ref] (us.south east) -- node[lbl, pos=.55] {created\_by} (pv.north west);
  \draw[own] (ct) -- node[lbl] {messages} (cm);
  \draw[ref] (ct.south) -- node[lbl, pos=.5, left] {company, role} (cx.north);
  \draw[own] (jo) -- node[lbl] {steps} (js);
  \draw[own] (pr) -- node[lbl, left] {versions} (pv);
  \draw[ref] ([xshift=6mm]pr.south) -- node[lbl, right] {active} ([xshift=6mm]pr.south |- pv.north);
  \draw[txt] (js) -- node[lbl, pos=.62] {prompt\_version\_id} (pv);
  \draw[txt] (lc) -- node[lbl] {job\_step\_id} (js);
  \draw[txt] (lc.north east) -- node[lbl, pos=.45] {prompt\_version\_id} (pv.south west);
  \draw[txt] (lc.west) -- node[lbl, pos=.4] {thread\_id} (ct.south east);
  \draw[txt] (lc) -- node[lbl] {priced by model name} (mp);
  \draw[ref] (rw.east) -- node[lbl, above] {job} (jo.west);
  \draw[ref] ([xshift=3mm]us.south) .. controls (92,-12) and (92,-48) .. node[lbl, pos=.3] {user} (lc.north east);
  \node[note, anchor=north, align=center] at (st.south) {no relations, but \code{mcp_servers.created_by}\\ points at \code{users} (emptied when the user goes)};
\end{tikzpicture}
\caption{The operational tables. Solid: a part deleted with its whole; dotted: a reference
  emptied when its target is deleted; dashed: an id held as text, with no foreign key. The two
  dashed boxes stand for tables of \cref{fig:data:law}.}\label{fig:data:ops}
\end{figure}

\section{The legal corpus and the register of obligations}\label{sec:data:corpus}

The application \code{corpus} holds the acts, their versions and their text cut into units and
chunks, and the obligations read out of them. An act is a row of \code{legal_documents}; each
edition fetched or uploaded is a row of \code{legal_document_versions}; the structure of a
version is a tree of \code{legal_units}; and the passages that retrieval searches are
\code{legal_chunks}. How a text gets there --- the sources, the parsers, the chunking --- is the
subject of \cref{ch:corpus}; how obligations are extracted and judged, of
\cref{ch:obligations}. On 13 November 2026 the tables held 78 acts, 71 versions, 28,886 units,
4,059 chunks and 5,960 obligations.

\begin{xltabular}{\textwidth}{@{}L{38mm}Y@{}}
\caption{\code{legal_documents}: one act or standard of the catalogue}\label{tab:data:ld}\\
\toprule
\hd{Column} & \hd{Meaning} \\
\midrule
\endfirsthead
\toprule
\hd{Column} & \hd{Meaning} \\
\midrule
\endhead
\bottomrule
\endlastfoot
\code{slug} & Unique; the key the catalogue (\file{corpus/catalog.yaml}) names the act by. \\
\code{title}, \code{title_en}, \code{short_title} & The official title, its English gloss and the short title that citations print (\enquote{MiCA}; for a standard its designation, \enquote{ISO/IEC 27001:2022}). \\
\code{jurisdiction} & \code{EU}, \code{LT} or \code{INTL} (an international standard). \\
\code{doc_type} & \code{regulation}, \code{directive}, \code{law}, \code{resolution}, \code{order}, \code{guideline}, \code{licensing}, \code{standard}, \code{other}. \\
\code{source_kind}, \code{external_id}, \code{source_url} & Where the text comes from: \code{cellar}, \code{etar}, \code{eseimas} (the legacy name of \code{etar}), \code{lb}, \code{url}, \code{upload}; the identifier at the source (a CELEX number, an e-TAR id) and the official link. \\
\code{language} & The language of the text, default \code{lt}. \\
\code{domains}, \code{applies_to} & JSON lists: the subject tags of the act and the institution types it typically governs (\code{casp}, \code{emi}, \dots{} or \code{all}) --- what the applicability screen filters on before a model reads the act. \\
\code{applicability_notes} & A note for that screen, written in the catalogue. \\
\code{status} & \code{in_force}, \code{not_yet_applicable}, \code{repealed}. \\
\code{enabled} & False hides the act from retrieval and analysis without deleting its versions. \\
\code{current_version} & The version in use; retrieval only ever reads chunks of the current version. \\
\code{last_checked_at}, \code{last_error} & The last refresh and its error. \\
\end{xltabular}

\begin{xltabular}{\textwidth}{@{}L{38mm}Y@{}}
\caption{\code{legal_document_versions}: one immutable edition of an act}\label{tab:data:lv}\\
\toprule
\hd{Column} & \hd{Meaning} \\
\midrule
\endfirsthead
\toprule
\hd{Column} & \hd{Meaning} \\
\midrule
\endhead
\bottomrule
\endlastfoot
\code{document}, \code{version_no} & The act and the edition's number; unique together. \\
\code{content_hash}, \code{label} & The hash of the raw bytes and what the source calls the edition (a consolidated CELEX such as \code{02023R1114-20240109}, an e-TAR edition id). A refresh that finds the same hash or label creates nothing, unless it is forced (the job's \code{force}, \file{corpus/ingest.py}). \\
\code{file_path}, \code{content_type}, \code{source_url} & The raw bytes on disk under the upload folder, so that a re-parse needs no re-fetch. \\
\code{effective_date} & The date the edition takes effect, when known. \\
\code{parser}, \code{text} & The structural parser that read it (empty for plain extraction) and the extracted plain text. \\
\code{char_count}, \code{unit_count}, \code{article_count}, \code{chunk_count}, \code{embed_model} & Counts and the embedding model of its chunks. \\
\code{status}, \code{error} & \code{pending}, \code{parsing}, \code{embedding}, \code{ready}, \code{failed}. \\
\code{obligations_extracted_at} & When obligation extraction last completed; empty means never, and a report then extracts first. \\
\code{scope}, \code{sanctions}, \code{changes} & JSON: a model's reading of the act's scope provisions and of its sanction provisions, and what the version changed against the one before (below). Empty means not read, or not recorded. \\
\end{xltabular}

The three JSON columns of a version have these shapes:

\begin{codeblock}
scope     = {types: [institution types, or "all"], reasoning,
             promptVersionId, model, readAt}
sanctions = {tiers: [{path, where, condition, documentId, infringedAct,
                      articles: [[from, to]], all, addressee, persons,
                      maxFineEur, maxFinePercent, benefitMultiple, measures}],
             note, provisions: [paths read], promptVersionId, model, readAt}
changes   = {from: previous version id, changed, added, removed: [paths]}
\end{codeblock}

\code{scope} is written by \file{corpus/scope.py} and lets the applicability screen keep an act
the catalogue's \code{applies_to} would leave out. In \code{sanctions}
(\file{corpus/sanctions.py}) each tier is a ceiling or a measure one provision lays down: its
\code{documentId} is the act it concerns (this one, another of the corpus, or empty with
\code{infringedAct} naming one the corpus does not hold), \code{articles} the article ranges
it reaches, each from its first to its last article by the numbers the act gives them ---
numbers alone in a reading made before 19~November 2026 --- (\code{all} when it reaches every
article), \code{persons} the kinds of person a tier binds when its provision binds only some
of those the act regulates --- kept on a tier without articles, and on one that names
articles beside another tier of the same articles for other persons (empty otherwise, and
absent from a reading made before that day) ---, \code{addressee} one of \code{legal},
\code{natural}, \code{any}, and \code{measures} codes from \code{withdrawal},
\code{restriction}, \code{cease}, \code{public_statement}, \code{management_ban},
\code{disgorgement}, \code{periodic_penalty}, \code{criminal}, \code{other}; at most
\code{MAX_TIERS = 24} are kept. \code{changes} is computed without a model when a new version
is activated (\code{version_changes} in \file{corpus/ingest.py}), article by article --- or by
top-level points in an act without articles --- comparing the texts with white space ignored.

\paragraph{Units and chunks.} A version's text is stored twice more: as a tree of units, which
is what a citation names, and as chunks, which is what retrieval searches
(\cref{tab:data:lu}).

\begin{xltabular}{\textwidth}{@{}L{38mm}Y@{}}
\caption{\code{legal_units} and \code{legal_chunks}}\label{tab:data:lu}\\
\toprule
\hd{Column} & \hd{Meaning} \\
\midrule
\endfirsthead
\toprule
\hd{Column} & \hd{Meaning} \\
\midrule
\endhead
\bottomrule
\endlastfoot
\multicolumn{2}{@{}l}{\textit{legal\_units}} \\
\code{version}, \code{parent}, \code{ordinal}, \code{depth} & The node's place in the version's tree, in document order (indexed by version and ordinal). \\
\code{unit_type} & \code{part}, \code{title}, \code{chapter}, \code{section}, \code{subsection}, \code{article}, \code{paragraph}, \code{point}, \code{subpoint}, \code{annex}, \code{recital}, \code{preamble}, \code{text}, \code{clause} (a standard's numbered clause, \enquote{6.1.2}). \\
\code{number}, \code{heading} & Its own number and heading; version, type and number together are indexed for the lookup of a cited article. \\
\code{path} & The citation it prints: \enquote{59 str. 1 d.}, \enquote{Art. 59(1)}. \\
\code{text} & Its \emph{full} text including its children, so that an article row alone answers \enquote{quote article 59}. \\
\midrule
\multicolumn{2}{@{}l}{\textit{legal\_chunks}} \\
\code{version}, \code{unit}, \code{ordinal} & The retrieval unit --- an article, or a run of paragraphs of a long one --- and its place (indexed by version and ordinal). \\
\code{document} & A copy of the act, so that search can filter by act without a join (indexed). \\
\code{path}, \code{heading}, \code{text}, \code{char_count} & What a search result shows; the GIN index covers \code{text}. \\
\code{embedding} & Under the HNSW index; empty until the embedding step has run, and search skips such rows. \\
\end{xltabular}

\begin{xltabular}{\textwidth}{@{}L{38mm}Y@{}}
\caption{\code{obligations}: one duty, or one requirement}\label{tab:data:ob}\\
\toprule
\hd{Column} & \hd{Meaning} \\
\midrule
\endfirsthead
\toprule
\hd{Column} & \hd{Meaning} \\
\midrule
\endhead
\bottomrule
\endlastfoot
\code{source_kind} & \code{legal} (from an act), \code{internal} (a company's own requirement, read from its policy), \code{standard} (a standard's requirement, extracted or entered by an administrator). \\
\code{document}, \code{version}, \code{unit} & The act, the version it was read from and the first provision it cites (for \code{legal} and \code{standard}). \\
\code{company}, \code{company_document} & For \code{internal}: the company and the policy it was read from. \\
\code{title}, \code{title_en}, \code{text}, \code{text_en} & The duty in Lithuanian and in English. \\
\code{addressee}, \code{conditions}, \code{deadline}, \code{evidence} & Who is bound, under what conditions, by when, and what would prove compliance. \\
\code{applies_to} & JSON list of institution types (or \code{all}) the duty binds; a report keeps the obligations that fit the company. \\
\code{category} & One of twenty: \code{licensing}, \code{registration}, \code{governance}, \code{policy}, \code{procedure}, \code{risk_management}, \code{ict_security}, \code{outsourcing}, \code{incident_reporting}, \code{reporting}, \code{disclosure}, \code{record_keeping}, \code{customer_protection}, \code{aml_kyc}, \code{capital}, \code{safeguarding}, \code{training}, \code{data_protection}, \code{ai_governance}, \code{other}. \\
\code{severity} & \code{critical}, \code{high}, \code{medium}, \code{low}. \\
\code{citations} & JSON: the provisions it rests on (below). \\
\code{status} & \code{proposed} (from the model), \code{approved}, \code{rejected}. A rejected row never reaches a report; an approved row survives re-extraction. \\
\code{prompt_version_id}, \code{model} & The prompt version and model that extracted it, as text. \\
\code{verification} & JSON: the judge's verdict (below); empty when never judged, or edited since. When none of the cited provisions is in the act's current version, code writes \code{unsupported} without a model call (\file{corpus/verify.py}). \\
\code{embedding}, \code{embedding_en} & Vectors of the title and text in each language. \\
\code{notes}, \code{reviewed_by}, \code{reviewed_at} & A reviewer's note, and who approved or rejected it. \\
\end{xltabular}

\begin{codeblock}
citations (legal, standard) = [{unitId, path}]
citations (internal)        = [{path: the section, documentId, documentTitle,
                                quote, start}]
verification = {verdict: supported | partly | unsupported,
                issue: none | duty | addressee | condition | deadline | merged,
                note, promptVersionId, model, checkedAt}
\end{codeblock}

For an own requirement, \code{start} is the character at which its quote stands in the
policy's text, or empty when the text does not hold it word for word. The verification
\enquote{orders a person's review and decides nothing} (\file{corpus/models.py}). Two indexes
serve the register: by act and status, and by company and status.

\section{Companies and their documents}\label{sec:data:companies}

The application \code{companies} holds the profile a report runs on, what the company holds and
has adopted, its roles, and its documents with their text, inventory and passages
(\cref{ch:company}). A company belongs to its owner; a document has an owner of its own, because
a file can be dropped into the chat before any company exists and attached later.

\begin{xltabular}{\textwidth}{@{}L{38mm}Y@{}}
\caption{\code{companies}: the profile}\label{tab:data:co}\\
\toprule
\hd{Column} & \hd{Meaning} \\
\midrule
\endfirsthead
\toprule
\hd{Column} & \hd{Meaning} \\
\midrule
\endhead
\bottomrule
\endlastfoot
\code{owner} & The user it belongs to (\code{PROTECT}). \\
\code{name}, \code{registration_code}, \code{country}, \code{legal_form} & Identity; country default \code{LT}. \\
\code{nace_codes}, \code{sector}, \code{activities} & JSON list of NACE codes; free text. \\
\code{size_category}, \code{employees}, \code{annual_turnover_eur} & Size: \code{micro}, \code{small}, \code{medium}, \code{large}, \code{unknown}. \\
\code{institution_type} & What the applicability screen filters on first: \code{bank}, \code{credit_union}, \code{payment_institution}, \code{emi}, \code{casp}, \code{crowdfunding}, \code{consumer_credit}, \code{real_estate_credit}, \code{currency_exchange}, \code{investment_firm}, \code{insurance}, \code{other_financial}, \code{none}, \code{unknown}. \\
\code{uses_ai}, \code{provides_ai}, \code{processes_personal_data}, \code{is_public_sector} & Tri-state flags: empty means nobody said, and the screen then keeps the acts the flag would gate rather than read silence as \enquote{no}. \code{ai_description} describes the AI use. \\
\code{financial_market} & JSON: the financial-market classification (market, service categories, authorisation). \\
\code{attributes} & JSON: an open bag of facts the conversation learns that have no column yet. \\
\end{xltabular}

\paragraph{What the company holds.} A row of \code{company_credentials} is a certificate,
licence, permit or registration (\code{kind}: \code{certificate}, \code{license},
\code{permit}, \code{registration}) with its issuer, identifier, validity and scope. Its
\code{status} is \code{claimed} until a person checks it against the issuer's register and sets
\code{verified} (with \code{verified_by}, \code{verified_at}); a claimed credential counts as
evidence only with that caveat, and an \code{evidence_document} may point at the document that
proves it. A row of \code{company_standards} is a standard of the corpus the company has adopted,
one per company and standard (unique), with a \code{status} of \code{certified},
\code{implementing} or \code{planned}, the scope it covers and the credential that certifies
it; a report that names no acts of its own assesses every adopted standard, and one that names
its acts assesses the standards among them, adopted or not (\code{adopted_standards} in
\file{analysis/pipeline.py}). A row of
\code{company_roles} is a role profile --- the compliance officer, for instance --- distilled
into \code{context}, the text a chat thread can adopt, with \code{sources} listing the document
sections it was built from (\code{[{documentId, title, section}]}); the name is unique within
the company.

\begin{xltabular}{\textwidth}{@{}L{38mm}Y@{}}
\caption{\code{company_documents}: one uploaded file}\label{tab:data:cd}\\
\toprule
\hd{Column} & \hd{Meaning} \\
\midrule
\endfirsthead
\toprule
\hd{Column} & \hd{Meaning} \\
\midrule
\endhead
\bottomrule
\endlastfoot
\code{owner}, \code{company} & The user who uploaded it (\code{PROTECT}) and the company it is attached to (empty until attached). \\
\code{title}, \code{kind} & \code{kind}: \code{policy}, \code{procedure}, \code{contract}, \code{register}, \code{other}. \\
\code{file_path}, \code{mime}, \code{size}, \code{sha256} & The file under the upload folder and its hash. \\
\code{text}, \code{char_count} & The extracted text --- what the model reads and what is chunked. \\
\code{extraction_status}, \code{extraction_error} & \code{pending}, \code{done}, \code{empty}, \code{failed}. \\
\code{inventory}, \code{inventory_key} & JSON: the evidence inventory, and the key it was computed for, \code{<sha256 of the text>:<prompt version id>:s<SCHEMA_VERSION>}; a new text, prompt version or inventory shape makes it stale (\file{companies/inventory.py}, \code{SCHEMA_VERSION = 2}). \\
\code{requirements_key} & \code{<sha256 of the file>:<prompt version id>} of the last reading of its own requirements; empty means never read. \\
\code{indexed_at} & When its chunks were last made. \\
\code{outline} & JSON: the headings found in the text, \code{[{title, level, offset}]}, so a reader can jump to \enquote{§7.2}. \\
\code{uploaded_at} & The moment of upload; with \code{owner}, the order of the owner's list (indexed). \\
\end{xltabular}

The inventory is the stored answer of the prompt \prompt{report.inventory}, whose JSON schema
fixes its keys:

\begin{codeblock}
inventory = {documentType: policy | procedure | register | contract | plan
                           | report | profile | other,
             summary, scope, controls: [{name, description, quote}],
             roles, registers, approval, reviewCycle, lawsReferenced,
             citations: [{act, article, paragraph, point, claim, quote}],
             weaknesses}
\end{codeblock}

\code{roles}, \code{registers}, \code{lawsReferenced} and \code{weaknesses} are lists of
short texts; \code{citations} are the references the citation audit checks
(\cref{ch:company}). A row of
\code{company_document_chunks} is a window of the text with its \code{ordinal},
\code{char_start}, the nearest \code{heading} above it (the section a citation of the chunk
names), \code{text} and \code{embedding}, under the HNSW and GIN indexes described above and
one by document and ordinal.

\section{Reports and the action register}\label{sec:data:analysis}

The application \code{analysis} holds the reports and everything a report is made of, and the
action register that outlives them. One report is one row of \code{reports} with its acts
(\code{report_acts}), findings (\code{report_findings}), actions (\code{report_actions}) and
citation-audit rows (\code{report_citations}); \code{company_actions} belongs to the company,
not to a report. How a report is made is the subject of \cref{ch:report}, how the register
lives on of \cref{ch:register}.

\begin{xltabular}{\textwidth}{@{}L{38mm}Y@{}}
\caption{\code{reports}: one gap analysis of one company}\label{tab:data:re}\\
\toprule
\hd{Column} & \hd{Meaning} \\
\midrule
\endfirsthead
\toprule
\hd{Column} & \hd{Meaning} \\
\midrule
\endhead
\bottomrule
\endlastfoot
\code{company}, \code{created_by}, \code{job} & The company, the user who started it, and the job of its latest run (\cref{ch:arch}); the job is replaced when the second part of a run is queued. \\
\code{status} & \code{queued}, \code{running}, \code{scope_review} (waiting for a person to approve the list of acts), \code{done}, \code{failed}, \code{cancelled}. \\
\code{language} & \code{lt} or \code{en}: the language the report is written in. \\
\code{scope} & JSON: what the run was asked to do (below). \\
\code{profile_snapshot} & JSON: the company as the run read it (below). \\
\code{title}, \code{summary} & The title and the executive summary (Markdown) the summary prompt wrote. \\
\code{stats} & JSON: counts, cost and provenance (below). \\
\code{error} & Why a run failed; \code{AI_BUDGET_EXHAUSTED: …} when the gateway refused for its budget. \\
\code{scope_approved_by}, \code{scope_approved_at} & Who approved the list of acts, and when. A run started without the review is stamped by the pipeline itself: a time and nobody. \\
\code{approved_by}, \code{approved_at}, \code{approval_note} & The sign-off; cleared again by any later review of a finding (a change of its override, note or cause). \\
\code{created_at}, \code{started_at}, \code{finished_at} & Its three moments; indexed by company and creation time. \\
\end{xltabular}

\begin{codeblock}
scope = {documentIds: [ids of legal_documents named by the user],
         review: true,     only when the scope review was asked for
         internal: true}   only when own requirements are to be assessed

stats = {acts: {applies, uncertain, not_applies,
                assessed, upcoming, unconfirmed, indirect},
         findings: {met, partial, missing, unknown, not_applicable},
         openBySeverity: {critical, high, medium, low},
         upcoming: {findings, openBySeverity},  only if an upcoming act has findings
         internal: {findings, openBySeverity},  only if own requirements have findings
         changes: {since, sinceDate, closed: {count, titles}, new, back},
                                                 only if an earlier report exists
         actions, reviewedObligations, tokens, costUsd,
         citations: {checked, match, mismatch, notFound, notInCorpus, unknown},
         provenance: {prompts: [{slug, version, model}], models,
                      corpus: [{act, versionNo, versionDate}]}}
\end{codeblock}

The \code{documentIds} of the scope are acts of the corpus, not company documents: an act
named there is assessed even when its verdict is \code{uncertain} or it binds the company only
indirectly (\code{assessment_of} in \file{analysis/pipeline.py}); one the catalogue screen
leaves out or the model judges not to apply is
not. \code{stats} is computed by \code{report_stats}
in \file{analysis/pipeline.py} when the run ends, and once before, when a run stops for the
scope review. The finding counts are by \emph{effective} status --- a reviewer's override
wins --- over the acts in force for the company; findings of acts assessed as preparation and
of the company's own requirements are counted apart, under \code{upcoming} and
\code{internal}. The counts are recomputed when a reviewer overrides a finding
(\code{finding_counts}); the rest is not. \code{tokens} and \code{costUsd} add up the model
calls of every job step filed under the report (\code{subject = report:<id>}), so both parts
of a run with a scope review are counted. \code{provenance} names each prompt version and
model that ran and, for every act the report judged rather than screened out, the version
number and date it read. \code{changes} (\code{changes} in \file{analysis/register.py})
compares each obligation with the last earlier finished report that assessed it: gaps
\code{closed}, gaps \code{new}, and obligations that went \code{back} from met to a gap, each
with a count and at most five titles.

\code{profile_snapshot} is the company as the API serves it with everything attached
(\code{serialize_company(…, with_related=True)} in \file{companies/serializers.py}): the
profile's columns under their wire names (\code{institutionType}, \code{usesAi},
\code{financialMarket}, \code{attributes}, \dots), and the lists \code{credentials},
\code{standards}, \code{roles} and \code{documents} --- the last with each document's
metadata, its inventory's \code{documentType} and one-line \code{inventorySummary}, never its
text. It is written by \code{freeze_profile} in \file{analysis/pipeline.py}: in the
\code{profile} step of the run's first part, and again when the second part starts. The second
part reads the company anew, so a report made with the scope review shows the company its
findings were judged on, with a credential verified or a document added while the report
waited. A report that runs without the review writes it once.

\begin{xltabular}{\textwidth}{@{}L{38mm}Y@{}}
\caption{\code{report_acts}: one act, as the report treated it}\label{tab:data:ra}\\
\toprule
\hd{Column} & \hd{Meaning} \\
\midrule
\endfirsthead
\toprule
\hd{Column} & \hd{Meaning} \\
\midrule
\endhead
\bottomrule
\endlastfoot
\code{document}, \code{version} & Links to the act and the version read (kept for navigation only). \\
\code{document_title}, \code{short_title}, \code{jurisdiction}, \code{doc_type}, \code{source_url} & Copies from the catalogue when the report ran. \code{doc_type} \code{standard} marks an adopted standard (applies by adoption, no model call); \code{internal} the row the company's own requirements hang on. \\
\code{applicability}, \code{confidence}, \code{reasoning}, \code{conditions_to_verify} & The model's verdict: \code{applies}, \code{not_applies} or \code{uncertain}; a confidence clamped to 0–1; its reason; the facts to confirm. \\
\code{applies_from} & The date the act binds the company from, when that is later than the day of the run. \\
\code{assessment} & What the report did with the act: \code{assessed}, \code{upcoming}, \code{unconfirmed}, \code{indirect}, \code{none}. Computed by code from the verdict; the scope review may change it, never the verdict. \\
\code{citations} & JSON: the scope provisions the verdict cites, \code{[{unitId, path, documentTitle}]}. \\
\code{screened_out}, \code{order} & True for an act the catalogue's screen left out without a model call; the order of listing. \\
\end{xltabular}

\begin{xltabular}{\textwidth}{@{}L{38mm}Y@{}}
\caption{\code{report_findings}: the verdict on one obligation}\label{tab:data:rf}\\
\toprule
\hd{Column} & \hd{Meaning} \\
\midrule
\endfirsthead
\toprule
\hd{Column} & \hd{Meaning} \\
\midrule
\endhead
\bottomrule
\endlastfoot
\code{act}, \code{obligation}, \code{source_kind} & Links to the report's act row and to the obligation; \code{legal}, \code{standard} or \code{internal}. \\
\code{document_title}, \code{short_title}, \code{obligation_title}, \code{obligation_text}, \code{category}, \code{deadline} & Copies of the act's titles and of the obligation (in the report's language when the obligation has an English text). \\
\code{obligation_reviewed} & Whether a person had approved the obligation when the report ran. \\
\code{status} & The model's verdict: \code{met}, \code{partial}, \code{missing}, \code{unknown} (not evidenced), \code{not_applicable}. \\
\code{severity} & The obligation's severity, copied. \\
\code{analysis}, \code{recommendation} & The model's assessment and recommendation. \\
\code{evidence}, \code{citations}, \code{consequence} & JSON (below). \\
\code{override_status}, \code{reviewer_note}, \code{reviewed_by}, \code{reviewed_at} & A reviewer's status and note; they win in every view and export. \\
\code{cause} & Why the gap exists, in the reviewer's words; only a person writes it. \\
\code{order} & Missing first, then partial, not evidenced, met, not applicable; graver first within each; the company's own requirements after the rest (indexed with \code{report}). \\
\end{xltabular}

\begin{codeblock}
evidence    = [{kind: "passage",    id: document id,   title, quote, start}
               {kind: "credential", id: credential id, title, quote}
               {kind: "profile",    id: company id,    title, quote}]
citations   = the obligation's citations, with documentTitle and url added
              (legal, standard); copied as they are (internal)
consequence = [{act, path, where, condition, maxFineEur, maxFinePercent,
                benefitMultiple, persons, measures, measuresFrom: [paths],
                via: {act, path}}]      (via: under a level 2 act only)
\end{codeblock}

A quote in \code{evidence} is cut to 600 characters, and \code{start} is the character at
which it stands in the document's text, or empty when the document does not hold it word for
word: the reader then sees what the model claimed and that it could not be placed. A report
made when passages were indexed chunks stored the chunk's ordinal (\code{chunk}) instead, and
the API resolves it when it serves the report (\file{analysis/serializers.py}).
\code{consequence} holds the sanction tiers that cover the provision as the corpus held them
when the report ran (\code{Consequences.of} in \file{corpus/sanctions.py}): the act's own, then
one other act's. It is empty for an act not yet binding the company, for a requirement of a
standard or of the company's own, for an act whose sanctions were never read, and for a
provision no tier reaches. Under a level~2 act the entries are those of the article of the
parent act it was adopted under, each with \code{via}, that act and article
(\cref{ch:obligations}); a report made before 16~November 2026 holds none for such an act.

\paragraph{Actions and citation-audit rows.} A row of \code{report_actions} is one step of the
plan (\cref{tab:data:rx}); a row of \code{report_citations} (\cref{tab:data:rc}) is one
reference a company document makes to a provision, with the audit's verdict. A reference
several documents make was judged once, and each document still has its own row.

\begin{xltabular}{\textwidth}{@{}L{38mm}Y@{}}
\caption{\code{report_actions}: one step of the action plan}\label{tab:data:rx}\\
\toprule
\hd{Column} & \hd{Meaning} \\
\midrule
\endfirsthead
\toprule
\hd{Column} & \hd{Meaning} \\
\midrule
\endhead
\bottomrule
\endlastfoot
\code{title}, \code{description}, \code{owner_role}, \code{category} & What to do, and the role that owns it. \\
\code{priority} & Its place in the plan, computed by code --- from what it closes (today's gaps before preparation for upcoming acts, then by severity and status), then its timeframe --- never asked of the model. \\
\code{effort}, \code{timeframe} & \code{low}, \code{medium}, \code{high}; \code{immediate}, \code{30_days}, \code{90_days}, \code{12_months}. \\
\code{finding_ids} & JSON list of the ids of the findings it closes; an action that closes no finding (correcting references, confirming facts) has an empty list. \\
\code{order} & The order of listing. \\
\end{xltabular}

\begin{xltabular}{\textwidth}{@{}L{38mm}Y@{}}
\caption{\code{report_citations}: one reference of a document, and the audit's verdict}\label{tab:data:rc}\\
\toprule
\hd{Column} & \hd{Meaning} \\
\midrule
\endfirsthead
\toprule
\hd{Column} & \hd{Meaning} \\
\midrule
\endhead
\bottomrule
\endlastfoot
\code{company_document}, \code{document_title} & The document that makes the reference (link and copy). \\
\code{act}, \code{article}, \code{paragraph}, \code{point}, \code{claim}, \code{quote} & The reference exactly as the document wrote it, and what it says the provision covers. \\
\code{status} & \code{found}, \code{not_found}, \code{act_not_in_corpus}. \\
\code{heading_match} & For a found provision: \code{match}, \code{mismatch}, \code{unknown}; empty otherwise. \\
\code{judged_by} & \code{model}, or \code{overlap} when the judge was unavailable and word overlap decided; empty when nothing was compared. \\
\code{legal_document}, \code{unit} & Links to the act and the provision found. \\
\code{short_title}, \code{unit_path}, \code{heading}, \code{note}, \code{excerpt}, \code{citation_url} & Copies of the provision found, the judge's note, an excerpt of the provision and its official link. \\
\code{order} & Problems first (mismatch, not found, unknown, act not in the corpus, match), then the documents' order. \\
\end{xltabular}

\begin{xltabular}{\textwidth}{@{}L{38mm}Y@{}}
\caption{\code{company_actions}: one action of the company's register}\label{tab:data:ca}\\
\toprule
\hd{Column} & \hd{Meaning} \\
\midrule
\endfirsthead
\toprule
\hd{Column} & \hd{Meaning} \\
\midrule
\endhead
\bottomrule
\endlastfoot
\code{company}, \code{report}, \code{plan_action_id} & The company; the signed-off report whose plan brought it in; the id of that plan's action. Unique by company and \code{plan_action_id}, so bringing a plan in twice adds nothing. \\
\code{title}, \code{description}, \code{category}, \code{priority}, \code{effort} & Copied from the plan's action. \\
\code{owner}, \code{due_date} & Who carries it out (first the plan's owner role) and by when (the sign-off date plus 14, 30, 90 or 365 days for the four timeframes); a person may change both. \\
\code{status} & \code{open}, \code{in_progress}, \code{done} (its owner's word), \code{closed}. \enquote{Late} is not stored: it is an open or in-progress action past its due date. \\
\code{items} & JSON: the open findings it is to close (below). \\
\code{note}, \code{closed_at}, \code{closed_by}, \code{closed_from}, \code{confirmed_report} & How it was closed: by a person with a reason in \code{note}, or by a later signed-off report (\code{confirmed_report}); \code{closed_from} is the state it was closed from, so a withdrawn sign-off can put it back. \\
\code{done_at} & When its owner last called it done: only a report made after that can send it back to work. \\
\end{xltabular}

\begin{codeblock}
items = [{obligationId, findingId, title, act, path, sourceKind, status,
          confirmed: null | {reportId, at}}]
\end{codeblock}

\section{Work products}\label{sec:data:programme}

The application \code{programme} has one table. A row of \code{work_products} is one version of
one product of a company --- a compliance policy, a structure with its process and RACI, or a
risk register (\cref{ch:programme}). Its \code{kind} is \code{policy}, \code{governance} or
\code{risks}, without a constraint; \code{version} counts up per company and kind (unique
together); \code{language} is \code{lt} or \code{en}; \code{status} is \code{queued},
\code{running}, \code{draft} (generated or being edited, waiting for a person), \code{approved}
(frozen: a change is a new version), \code{failed} or \code{cancelled}. Beside the
\code{job}, \code{title} and \code{error} it carries four JSON columns and three stamps of who
did what: \code{created_by}, \code{edited_by} with \code{edited_at}, and \code{approved_by}
with \code{approved_at} and \code{approval_note}.

\begin{codeblock}
content (policy)     = {title,
                        parts: [{key, heading, text, decisions, sources,
                                 legalBasis}],
                        hierarchy: [{documentId, document, level, topic, note}],
                        missing: [{level, title, purpose, supports, priority,
                                   legalBasis, partlyCoveredBy}]}
content (governance) = {title, positions, gaps, structure, process, raci}
content (risks)      = {title, criteria: {likelihood, impact, calibration},
                        appetite, scenarios}
checks = [{rule, ok, blocking, detail}]
basis  = {reportId, products: {<kind>: {id, version, status}},
          documents: [{id, title, sha256}]}
stats  = {tokens, costUsd, provenance: {prompts: [{slug, version, model}],
                                        models}}
\end{codeblock}

\code{content} is the product in the stored shape its kind's \code{validate} gives it
(\file{programme/policy.py}, \file{governance.py}, \file{risks.py}); \code{checks} are the
rules of the kind as the content last met or broke them, and a broken \code{blocking} rule
keeps the product from being approved; \code{basis} names what the run read --- the report,
the earlier products, the documents with their hashes --- so that a reader can tell what the
product rests on.

\section{The tables that run the system}\label{sec:data:ops}

The remaining tables, drawn in \cref{fig:data:ops}, are accounts and sessions, prompts, the
queue, the record of model calls, the chat and the audit trail.


\paragraph{Users and sessions.} A row of \code{users} is an account, and every other table that
names a person points here.

\begin{xltabular}{\textwidth}{@{}L{38mm}Y@{}}
\caption{\code{users} and \code{sessions}}\label{tab:data:us}\\
\toprule
\hd{Column} & \hd{Meaning} \\
\midrule
\endfirsthead
\toprule
\hd{Column} & \hd{Meaning} \\
\midrule
\endhead
\bottomrule
\endlastfoot
\multicolumn{2}{@{}l}{\textit{users}} \\
\code{email}, \code{display_name} & The e-mail is unique, and unique ignoring case as well (a functional index on its lower case). \\
\code{password_hash} & The hash of the password; never the password. \\
\code{role} & \code{admin} or \code{user}. \\
\code{active} & The kill switch every authenticated request checks again; the plain delete of an account deactivates it, and only an administrator's purge deletes it with its data (\cref{tab:data:ondelete}). \\
\code{daily_budget_usd} & What the account may spend on model calls per UTC day: empty means the installation's default (\code{LLM_DAILY_BUDGET_USD}), 0 means unlimited (\cref{ch:models}). \\
\code{created_at}, \code{updated_at}, \code{last_login_at} & \\
\midrule
\multicolumn{2}{@{}l}{\textit{sessions}} \\
\code{user}, \code{token_hash} & The account, and the SHA-256 of the bearer token the client holds (unique); the token itself is never stored. \\
\code{expires_at}, \code{created_at}, \code{user_agent}, \code{ip} & Expiry and where the session came from; indexed by user and expiry. At sign-in only the newest \code{SESSIONS_PER_USER = 10} of an account are kept (\file{users/auth.py}). \\
\end{xltabular}

\paragraph{Prompts.} The instructions a model is given are stored as prompts with immutable
versions (\cref{tab:data:pr}); editing a prompt adds a version, and activating one moves
\code{active_version}. How shipped defaults become versions is the subject of
\cref{ch:prompts}. On 13 November 2026 the tables held 23 prompts and 85 versions, and every
active version was a shipped default; a new default of the chat's prompt was seeded on 16 and
again on 17 November 2026, new defaults of the sanctions prompt (twice) and of the
inventory prompt on 19 November and one more of the sanctions prompt on 20 November, so
they hold 91 versions, and every active version is still
a shipped default.

\begin{xltabular}{\textwidth}{@{}L{38mm}Y@{}}
\caption{\code{prompts} and \code{prompt_versions}}\label{tab:data:pr}\\
\toprule
\hd{Column} & \hd{Meaning} \\
\midrule
\endfirsthead
\toprule
\hd{Column} & \hd{Meaning} \\
\midrule
\endhead
\bottomrule
\endlastfoot
\multicolumn{2}{@{}l}{\textit{prompts}} \\
\code{slug}, \code{name}, \code{description} & The unique name the code calls it by (\prompt{report.gap_analysis}). \\
\code{kind} & \code{chat}, \code{pipeline} or \code{corpus}. \\
\code{variables}, \code{output_schema} & JSON: the placeholders it may use, and the JSON schema its answer must match. \\
\code{active_version} & The version in use. \\
\midrule
\multicolumn{2}{@{}l}{\textit{prompt\_versions}} \\
\code{prompt}, \code{version_no} & Unique together. \\
\code{text}, \code{model}, \code{temperature}, \code{max_tokens} & The text and the model settings; \code{model} is a string of the form \code{provider:model}. \\
\code{source}, \code{default_hash} & \code{default} for a text shipped in \file{prompts/defaults/*.yaml}, whose hash names the exact shipped text; \code{user} for one written in the editor. \\
\code{note}, \code{created_by}, \code{created_at} & \\
\end{xltabular}

\paragraph{Jobs and steps.} The queue is two tables, and \cref{ch:arch} describes how they
behave as one: claiming, locking, heartbeats, retries and cancelling. \Cref{tab:data:jobs} lists
their columns; the JSON among them have these shapes:

\begin{codeblock}
jobs.progress   = {step, percent, message}
jobs.result     = what the handler returned; report.analyze:
                  {reportId, findings, actions} or {reportId, waiting: "scope"};
                  programme.generate: {productId, checks: how many broken}
job_steps.input_ref, job_steps.output
                = what the step was given and what it counted; the
                  gap_analysis output, for instance: {met, partial, missing,
                  unknown, not_applicable, documents: "all" | "nearest",
                  quotesPlaced, quotesUnplaced}
job_steps.usage = {calls, inputTokens, outputTokens, costUsd,
                   models: [{model, calls, inputTokens, outputTokens, costUsd}]}
\end{codeblock}

\begin{xltabular}{\textwidth}{@{}L{38mm}Y@{}}
\caption{\code{jobs} and \code{job_steps}}\label{tab:data:jobs}\\
\toprule
\hd{Column} & \hd{Meaning} \\
\midrule
\endfirsthead
\toprule
\hd{Column} & \hd{Meaning} \\
\midrule
\endhead
\bottomrule
\endlastfoot
\multicolumn{2}{@{}l}{\textit{jobs}} \\
\code{kind}, \code{payload}, \code{subject} & The handler, its JSON input and what the job is about (\code{report:<id>}; indexed). \\
\code{status} & \code{queued}, \code{running}, \code{done}, \code{failed}, \code{cancelled}. \\
\code{priority}, \code{run_after} & The order of claiming and the earliest moment; one index serves the claim (status, \code{run_after}, priority, creation). \\
\code{attempts}, \code{max_attempts}, \code{cancel_requested}, \code{locked_by}, \code{heartbeat_at} & Retries, cancelling and the lock of the worker that holds it. \\
\code{progress}, \code{result}, \code{error} & JSON progress and result; the error text. \\
\code{created_by}, \code{created_at}, \code{started_at}, \code{finished_at} & Who caused it (empty for the scheduler's) and its moments. \\
\midrule
\multicolumn{2}{@{}l}{\textit{job\_steps}} \\
\code{job}, \code{ordinal}, \code{name} & The step's place in its job (indexed): \code{profile}, \code{documents}, \code{citations}, \code{applicability}, … for a report. \\
\code{status} & \code{running}, \code{done}, \code{failed}, \code{skipped}. \\
\code{prompt_version_id}, \code{model} & The prompt version and model of its calls, as text: the step's record outlives the version row. \\
\code{input_ref}, \code{output}, \code{usage} & JSON, as above; \code{usage} is read back from the step's rows of \code{llm_calls} when the step ends. \\
\code{duration_ms}, \code{error}, \code{started_at}, \code{finished_at} & \\
\end{xltabular}

\paragraph{The record of model calls.} A row of \code{llm_calls} is one call of a model, and the
table is what the usage page, the daily budget and every cost figure of this book are computed
from (\cref{ch:models,ch:control}).

\begin{xltabular}{\textwidth}{@{}L{38mm}Y@{}}
\caption{\code{llm_calls}: one model call}\label{tab:data:lc}\\
\toprule
\hd{Column} & \hd{Meaning} \\
\midrule
\endfirsthead
\toprule
\hd{Column} & \hd{Meaning} \\
\midrule
\endhead
\bottomrule
\endlastfoot
\code{kind}, \code{purpose} & \code{chat}, \code{job_step}, \code{embed}, \code{prompt_test}; a short free label (at most 80 characters). \\
\code{user} & The account charged; empty for work charged to nobody, such as corpus extraction inside a report. \\
\code{thread_id}, \code{job_step_id}, \code{prompt_version_id} & What the call was for, as text ids with no foreign key; \code{job_step_id} is indexed, since a step's usage and a report's cost are read through it. \\
\code{model} & The model as the gateway names it. \\
\code{input_tokens}, \code{cached_input_tokens}, \code{cache_write_tokens}, \code{output_tokens}, \code{reasoning_tokens} & Token counts; of the input, those read from the provider's prompt cache and those written to it, which are billed at other prices. \\
\code{cost_usd} & The cost in US dollars, six decimals, computed from the gateway's price list when the row is written. \\
\code{duration_ms}, \code{tool_calls}, \code{outcome} & Duration, the JSON list of tools a chat turn called, and the outcome (\code{ok} by default; free text up to 40 characters). \\
\code{transcript} & JSON for a chat turn: \code{{user, assistant}}, the text of the question and of the answer --- the audit of what the assistant told whom. Erased when the account is purged (\file{users/purge.py}); the row itself stays, unlinked. \\
\code{created_at} & Indexed alone and with \code{user}. \\
\end{xltabular}

Two small tables of the same application stand beside it. \code{model_prices}, keyed by the
model's id, keeps the last known US dollars per token of each model --- input, output, and,
when the gateway names them, input read from and written to the cache --- refreshed from the
gateway's list and used when the list cannot be reached. \code{mcp_servers} holds an external
MCP server whose tools the chat assistant may call: a unique \code{name} (the prefix of its
tools), \code{url}, \code{transport} (\code{http} or \code{sse}), \code{headers} with its
credentials, \code{allowed_tools} (empty means every tool), \code{enabled}, and the last
successful listing of its tools (\code{tools}, \code{[{name, description}]}).

\paragraph{The chat.} A conversation is a thread of one user with its messages
(\cref{tab:data:ct}). The console appends the user's messages and the assistant service its
replies, both through \file{chat/store.py}, which refuses a message whose content does not
carry exactly one role and a stored message that would change its role (\cref{ch:chat}).

\begin{xltabular}{\textwidth}{@{}L{38mm}Y@{}}
\caption{\code{chat_threads} and \code{chat_messages}}\label{tab:data:ct}\\
\toprule
\hd{Column} & \hd{Meaning} \\
\midrule
\endfirsthead
\toprule
\hd{Column} & \hd{Meaning} \\
\midrule
\endhead
\bottomrule
\endlastfoot
\multicolumn{2}{@{}l}{\textit{chat\_threads}} \\
\code{user} & The one user the conversation belongs to. \\
\code{company}, \code{role} & Optionally the company it is about and the role (\code{company_roles}) whose context it adopts. \\
\code{title} & Named by the first message, by a model, or by the user. \\
\code{model} & The gateway model the user picked (\code{gateway:azure/…}); empty means the chat prompt's own. \\
\code{message_count}, \code{archived}, \code{created_at}, \code{updated_at} & Indexed by user, archived and last update. \\
\midrule
\multicolumn{2}{@{}l}{\textit{chat\_messages}} \\
\code{id} & The id the chat client gave the message (see \cref{sec:data:conventions}). \\
\code{thread}, \code{parent_id} & The thread, and the message it follows, as text. \\
\code{format} & The format name the writer gives; the console and the assistant service write \code{ai-sdk/v6}. \\
\code{role} & The message's one role, read out of its content when stored. \\
\code{content} & JSON: the message as the chat client holds it, at most \code{MAX_CONTENT_CHARS = 400000} characters. \\
\code{created_at} & Kept when a message is sent again, so history never reorders (indexed with the thread). \\
\end{xltabular}

\paragraph{The audit trail and the rate limiter.} A row of \code{audit_log} records one
mutation an administrator may need to see --- a sign-in, a change of a user, a prompt
activation, a change of the corpus, a report's sign-off (\cref{tab:data:au}). It is written
only through \code{audit()} in \file{audit/log.py}. A row of \code{rate_limit_windows} is one
fixed window of a limiter: its \code{key}, the \code{count} of requests in it and
\code{reset_at}, shared by every web worker of the API (\file{common/ratelimit.py},
\cref{ch:security}).

\begin{xltabular}{\textwidth}{@{}L{38mm}Y@{}}
\caption{\code{audit_log}: one recorded action}\label{tab:data:au}\\
\toprule
\hd{Column} & \hd{Meaning} \\
\midrule
\endfirsthead
\toprule
\hd{Column} & \hd{Meaning} \\
\midrule
\endhead
\bottomrule
\endlastfoot
\code{user} & Who did it; empty for an action without a signed-in user. \\
\code{action} & A dotted name: \code{reports.approve}, \code{reports.approve_scope}, \code{corpus.delete}, \dots \\
\code{target_type}, \code{target_id} & What it was done to, as text (indexed together). \\
\code{details} & JSON: what changed, as the action records it --- for a scope approval, each act whose treatment the person changed. \\
\code{ip}, \code{created_at} & Where from, and when (indexed newest first). \\
\end{xltabular}

\section{What a report copies, and what it does not}\label{sec:data:snapshots}

A compliance report is evidence of what was found on a given day, and it must read the same
later, when the corpus has a new version of an act, an obligation has been reworded, a
document has been replaced or the company's profile has changed. The model files state the
rule: a report \enquote{snapshots every title and text it shows (act titles, obligation texts,
citation paths, provision headings), so it reads the same after the corpus, the obligations,
the documents or the profile change; the foreign keys are kept only for navigation}
(\file{analysis/models.py}).

\begin{codebox}{What is copied into a report when it runs}
\begin{itemize}
  \item The company: \code{reports.profile_snapshot}, the whole profile with its credentials,
    adopted standards, roles and the list of documents with their inventory summaries. A
    report that stopped for the scope review copies it again when its second part starts.
  \item Every act: its title, short title, jurisdiction, type and link
    (\code{report_acts}), and the version read (a link, and its number and date in
    \code{stats.provenance.corpus}).
  \item Every obligation judged: its title and text in the report's language, its category,
    severity and deadline, whether a person had approved it, and the provisions it cites with
    their paths and official links (\code{report_findings}).
  \item What a breach can cost: the sanction tiers that covered the provision on the day
    (\code{report_findings.consequence}).
  \item The evidence: each quote the model gave, with the place it stands in its document on
    the day (\code{start}).
  \item The citation audit: each reference as written, the provision found, its heading, an
    excerpt and its link (\code{report_citations}).
  \item Which prompt version and model made each step, and what the run cost
    (\code{stats.provenance}, \code{tokens}, \code{costUsd}).
\end{itemize}
\end{codebox}

What is not copied is as deliberate. The full text of an act and of a company document stays
in its own table: a report holds quotes, excerpts and offsets into them, not the texts. A quote's
\code{start} is an offset into the document's text as it was; if the document is deleted the
link is emptied, and the quote remains. The prompts are named, not copied: provenance gives the
slug, the version number and the model, and the version itself stays, since prompt versions
are immutable and no route deletes one. The job's steps, which hold the per-step record, are not purged while the report
points at the job (\cref{ch:arch}). Two things in a served report are computed afresh on
purpose: \code{outdatedActs}, the acts that have received a newer version since the report read
them (so that a reader knows to run it again), and each finding's \code{history}, how its
obligation stood in the company's earlier reports (\file{analysis/serializers.py}).

The action register copies in the same way: an item of \code{company_actions.items} holds the
finding's title, act and provision path as they were at the sign-off, and names the obligation
by id. An obligation that was never approved is replaced when its act is read again, so its id
may disappear; the API then serves the item without its obligation, as one that no later
report can confirm (\code{standing} in \file{analysis/register.py}).

\section{Who writes which columns}\label{sec:data:writers}

The tables mix three kinds of content, and the line between them is the line this book draws
everywhere: what a model wrote, what the code computed, and what a person decided. Code stores
a pipeline prompt's answer only after it has been held to the prompt's strict JSON schema, its
values checked against the allowed ones, and every reference the model gave mapped (a provision \code{P3}, an obligation \code{O7}, a document \code{D4}) back to a real
row; a reference that maps to nothing is dropped.

\begin{aibox}{Columns that hold what a model answered}
\code{legal_document_versions.scope} and \code{.sanctions}; the text fields of an extracted
\code{obligations} row (title, text, addressee, conditions, deadline, evidence, category,
severity, \code{applies_to}) and its \code{verification}; \code{company_documents.inventory};
in a report, \code{report_acts.applicability}, \code{confidence}, \code{reasoning},
\code{conditions_to_verify} and \code{applies_from}, \code{report_findings.status},
\code{analysis}, \code{recommendation} and the quotes of \code{evidence},
\code{report_citations.heading_match} and \code{note} when \code{judged_by} is \code{model},
the text of \code{report_actions} and which drafted actions merge into one
(\prompt{report.plan_merge}), and \code{reports.title} and \code{summary};
\code{work_products.content} before a person edits it; \code{chat_threads.title} when a model
named the thread, and the assistant's messages.
\end{aibox}

\begin{codebox}{Columns the code computes}
\code{report_acts.assessment} (from the verdict, the act's type and the report's scope);
\code{report_findings.consequence}, \code{order} and the \code{start} of each quote;
\code{report_actions.priority} and \code{finding_ids} (a merged action closes what its drafts
closed, and a finding its planning call left without an action joins the nearest one,
\code{_nearest_action} in \file{analysis/pipeline.py}); \code{reports.stats};
\code{legal_document_versions.changes}; \code{work_products.checks}; \code{company_actions}
items, due dates and their confirmation by a later report; every count, key, hash and cost.
\end{codebox}

\begin{humanbox}{Columns that record a person's decision}
\code{obligations.status} (approving or rejecting), \code{reviewed_by} and \code{notes};
\code{company_credentials.status} set to \code{verified}; the scope review
(\code{reports.scope_approved_by}, and the \code{assessment} of the acts it changes); a
finding's \code{override_status}, \code{reviewer_note} and \code{cause}; the sign-off
(\code{reports.approved_by}, \code{approval_note}); the owner, due date, state and closing note
of an action in the register; the edit and approval of a work product (\code{edited_by},
\code{approved_by}); a prompt version of source \code{user}. No pipeline sets any of them.
Three can also be set through the chat, where the model calls a tool in the user's name and
with the user's rights (\cref{ch:chat}): approving a report's scope (\code{approveReportScope},
after the user has accepted the list and its price), overriding a finding
(\code{overrideFinding}, with the user's reason) and marking a credential verified
(\code{updateCredential}, or \code{addCredential} with that status, on the user's word that the
credential was checked). The row then names the user, as if the console had been used. No
tool writes a finding's \code{cause}.
\end{humanbox}

\section{How the schema came to be}\label{sec:data:history}

The tables are made and changed by Django migrations,
one folder per application. Their names are the short history of the model
(\cref{tab:data:migrations}). On 13 November 2026 the live database stood at the last migration
listed here for \code{analysis}, \code{assistant}, \code{corpus}, \code{companies}, \code{programme},
\code{prompts}, \code{users}, \code{chat}, \code{audit}, \code{jobs} and \code{ops}.

\begin{xltabular}{\textwidth}{@{}L{20mm}Y@{}}
\caption{The migrations of each application, by file name}\label{tab:data:migrations}\\
\toprule
\hd{Application} & \hd{Migrations after \code{0001_initial}} \\
\midrule
\endfirsthead
\toprule
\hd{Application} & \hd{Migrations after \code{0001_initial}} \\
\midrule
\endhead
\bottomrule
\endlastfoot
\code{corpus} & 0002 \code{obligations_extracted_at} · 0003 source kinds changed · 0004 \code{obligation_embedding_en} · 0005 \code{standards} · 0006 version \code{scope} · 0007 obligation \code{verification} · 0008 version \code{sanctions} · 0009 version \code{changes} \\
\code{companies} & 0002 document \code{indexed_at}, \code{inventory} and more · 0003 \code{companyrole}, credential evidence document and more · 0004 \code{companystandard} · 0005 document \code{requirements_key} \\
\code{analysis} & 0002 \code{initial} · 0003 \code{finding_obligation_reviewed} · 0004 report approval note and time · 0005 \code{recount_effective_findings} · 0006 \code{reportcitation} · 0007 act \code{doc_type} · 0008 act \code{assessment} · 0009 report \code{scope_review} · 0010 finding \code{consequence}, \code{cause} · 0011 \code{companyaction} \\
\code{assistant} & 0002 call cost and purpose · 0003 call kinds changed · 0004 \code{mcpserver} · 0005 \code{modelprice} · 0006 call job-step index · 0007 \code{prompt_cache_accounting} \\
\code{users} & 0002 user token budget · 0003 token budget replaced by \code{daily_budget_usd} \\
\code{chat} & 0002 \code{initial} · 0003 thread \code{role} · 0004 thread \code{model} \\
\code{prompts} & 0002 version \code{default_hash}, \code{source} \\
\code{programme}, \code{jobs}, \code{audit}, \code{ops} & none \\
\end{xltabular}
